% Options for packages loaded elsewhere
\PassOptionsToPackage{unicode}{hyperref}
\PassOptionsToPackage{hyphens}{url}
%
\documentclass[
]{article}
\usepackage{amsmath,amssymb}
\usepackage{iftex}
\ifPDFTeX
  \usepackage[T1]{fontenc}
  \usepackage[utf8]{inputenc}
  \usepackage{textcomp} % provide euro and other symbols
\else % if luatex or xetex
  \usepackage{unicode-math} % this also loads fontspec
  \defaultfontfeatures{Scale=MatchLowercase}
  \defaultfontfeatures[\rmfamily]{Ligatures=TeX,Scale=1}
\fi
\usepackage{lmodern}
\ifPDFTeX\else
  % xetex/luatex font selection
\fi
% Use upquote if available, for straight quotes in verbatim environments
\IfFileExists{upquote.sty}{\usepackage{upquote}}{}
\IfFileExists{microtype.sty}{% use microtype if available
  \usepackage[]{microtype}
  \UseMicrotypeSet[protrusion]{basicmath} % disable protrusion for tt fonts
}{}
\makeatletter
\@ifundefined{KOMAClassName}{% if non-KOMA class
  \IfFileExists{parskip.sty}{%
    \usepackage{parskip}
  }{% else
    \setlength{\parindent}{0pt}
    \setlength{\parskip}{6pt plus 2pt minus 1pt}}
}{% if KOMA class
  \KOMAoptions{parskip=half}}
\makeatother
\usepackage{xcolor}
\setlength{\emergencystretch}{3em} % prevent overfull lines
\providecommand{\tightlist}{%
  \setlength{\itemsep}{0pt}\setlength{\parskip}{0pt}}
\setcounter{secnumdepth}{-\maxdimen} % remove section numbering
\ifLuaTeX
  \usepackage{selnolig}  % disable illegal ligatures
\fi
\IfFileExists{bookmark.sty}{\usepackage{bookmark}}{\usepackage{hyperref}}
\IfFileExists{xurl.sty}{\usepackage{xurl}}{} % add URL line breaks if available
\urlstyle{same}
\hypersetup{
  hidelinks,
  pdfcreator={LaTeX via pandoc}}

\author{}
\date{}

\begin{document}

\hypertarget{satspath-estado-actual-p2p-y-roadmap-v0.2}{%
\section{SatsPath: Estado Actual (P2P) y Roadmap
v0.2}\label{satspath-estado-actual-p2p-y-roadmap-v0.2}}

Este documento resume el estado actual del proyecto (tras alcanzar el
100\% de cumplimiento de la especificación \textbf{SatsPath Protocol
v0.1}) y detalla las funcionalidades que el PDF original especifica como
pendientes para futuras versiones (``Future behavior'' / ``Production''
/ ``Do not build yet'').

\begin{center}\rule{0.5\linewidth}{0.5pt}\end{center}

\hypertarget{lo-que-ya-funciona-v0.1-completada-p2p}{%
\subsection{1. Lo que ya funciona (v0.1 Completada +
P2P)}\label{lo-que-ya-funciona-v0.1-completada-p2p}}

Hemos superado los requisitos básicos del Hackathon. Actualmente el
sistema cuenta con:

\begin{itemize}
\tightlist
\item
  \textbf{100\% Cumplimiento de la Spec v0.1}: Resolver, Router,
  perfiles firmados, rotación de llaves, verificación de firmas y
  estimación de comisiones.
\item
  \textbf{Integración P2P Funcional}: Los perfiles pueden publicarse y
  resolverse de manera descentralizada utilizando Holepunch / Pear,
  eliminando la necesidad de un servidor central
  (\texttt{satspath\ wallet\ publish}).
\item
  \textbf{Resolución Multi-backend}: Además del registro local y P2P, se
  soportan resolutores HTTP (\texttt{.well-known}) y Nostr (NIP-05).
\item
  \textbf{CLI y Daemon Nativos y Dockerizados}: Entorno seguro y
  funcional tanto vía Docker como compilado en Rust de forma nativa
  (\texttt{satspath-cli}, \texttt{satspathd}).
\item
  \textbf{Pruebas de Propiedad (Ownership Proofs)}: Sistema avanzado de
  mitigación de ataques que supera la especificación original.
\end{itemize}

\begin{center}\rule{0.5\linewidth}{0.5pt}\end{center}

\hypertarget{funcionalidades-faltantes-roadmap-v0.2-producciuxf3n}{%
\subsection{2. Funcionalidades Faltantes (Roadmap v0.2 /
Producción)}\label{funcionalidades-faltantes-roadmap-v0.2-producciuxf3n}}

Al revisar el documento \texttt{SatsPath\ Protocol\ v0.1.pdf}, las
siguientes funcionalidades fueron explícitamente marcadas como
\textbf{``Future behavior''}, \textbf{``Production rule''} o
\textbf{``Do not build yet''} (Sección 34). Esto es lo que falta
construir para llevar el protocolo a producción real:

\hypertarget{a.-ejecuciuxf3n-real-de-pagos-payment-execution}{%
\subsubsection{A. Ejecución Real de Pagos (Payment
Execution)}\label{a.-ejecuciuxf3n-real-de-pagos-payment-execution}}

Actualmente el sistema genera \textbf{Intenciones de Pago} (URIs,
códigos QR, facturas LN) y simula el éxito. Falta integrar la ejecución
real:

\begin{itemize}
\tightlist
\item
  \textbf{Lightning Network Real (§18.1 y §34)}: Pagar facturas
  automáticamente desde un nodo Lightning integrado.
\item
  \textbf{Transmisión On-chain Real (§34)}: Construcción de PSBTs
  (Partially Signed Bitcoin Transactions), firma de transacciones y
  \emph{broadcasting} a la mempool.
\item
  \textbf{Implementación Real de Ark (§18.3 y §34)}: Conectarse a un ASP
  (Ark Service Provider) real para transferir VTXOs.
\item
  \textbf{Ejecución de Pagos Divididos / Split Payments (§24)}: La
  estructura de datos ya está diseñada, pero el router aún no soporta la
  resolución recursiva y ejecución por lotes (\emph{batching}) para
  múltiples receptores.
\item
  \textbf{Custodial Escrow (§34)}: Depósito de garantía en custodia
  temporal (marcado como ``Do not build yet'').
\end{itemize}

\hypertarget{b.-privacidad-y-criptografuxeda}{%
\subsubsection{B. Privacidad y
Criptografía}\label{b.-privacidad-y-criptografuxeda}}

\begin{itemize}
\tightlist
\item
  \textbf{Silent Payments (§13 y §18.2)}: Para producción, la
  especificación exige no reutilizar direcciones estáticas. Se debe
  implementar generación de direcciones frescas o \textbf{Silent
  Payments} (BIP-352) para privacidad On-chain.
\item
  \textbf{Cálculo Real de Tamaño de Transacción (§18.2)}: Actualmente se
  usan estimaciones estáticas (ej. 140 vB para P2WPKH). En producción,
  se debe armar el PSBT para calcular los \emph{virtual bytes} exactos.
\item
  \textbf{Campos de Perfil Encriptados (§26 y §27)}: Para evitar fugas
  de privacidad en registros públicos, la especificación sugiere
  encriptación de campos usando secretos compartidos (ECDH).
\end{itemize}

\hypertarget{c.-seguridad-avanzada-y-modelo-de-amenazas-26}{%
\subsubsection{C. Seguridad Avanzada y Modelo de Amenazas
(§26)}\label{c.-seguridad-avanzada-y-modelo-de-amenazas-26}}

Las mitigaciones a futuro para vectores de ataque avanzados incluyen:

\begin{itemize}
\tightlist
\item
  \textbf{DNSSEC / BIP-353}: Validación criptográfica de resolución de
  nombres a través de DNS para mitigar registros de alias falsos.
\item
  \textbf{Logs de Transparencia (Transparency Logs)}: Para auditar y
  evitar la manipulación de registros en servidores centrales/HTTP.
\item
  \textbf{Políticas de Recuperación de Llaves (Multisig)}: Para mitigar
  la pérdida de la llave de identidad.
\item
  \textbf{Invitaciones Firmadas y Verificación de Dominio (§27)}: Para
  mitigar ataques de enlaces de invitación maliciosos.
\item
  \textbf{Validación de Pruebas y Reputación de Servidores Ark (§27)}:
  Para proteger a los usuarios de ASPs maliciosos.
\end{itemize}

\hypertarget{d.-experiencia-de-usuario-client-ux}{%
\subsubsection{D. Experiencia de Usuario (Client /
UX)}\label{d.-experiencia-de-usuario-client-ux}}

\begin{itemize}
\tightlist
\item
  \textbf{Billetera Móvil Completa (§34)}: El entregable actual es un
  Daemon y un CLI. Se requiere una aplicación móvil nativa (iOS/Android)
  para usuarios finales (``Do not build yet'').
\item
  \textbf{Soporte Completo para BOLT12 (§18.1)}: Obtención e interacción
  nativa con Ofertas BOLT12.
\end{itemize}

\begin{center}\rule{0.5\linewidth}{0.5pt}\end{center}

\hypertarget{conclusiuxf3n}{%
\subsection{Conclusión}\label{conclusiuxf3n}}

El entregable actual es un \textbf{prototipo funcional} perfecto para el
MVP y el hackathon, cumpliendo con la filosofía de \emph{``Trust
cryptographic signatures, not servers''}. El próximo gran salto (v0.2)
será integrar los motores de firmas para mover fondos reales y añadir
\emph{Silent Payments} para garantizar la privacidad.

\end{document}
